% ============================================================
% main.tex — selve handout-en
%
% Denne fila er struktur og innhold. Alle typografiske
% valg bor i handoutstyle.sty, som også laster babel.
%
% Kompiler med pdfLaTeX:
%   latexmk main
% ============================================================


\documentclass[11pt]{article}


% ============================================================
% LAST INN HANDOUT-STIL
%
% Med [unummerert] mister seksjonene nummeret sitt, og
% oppgaver og teoremer telles flatt gjennom dokumentet:
% Oppgave 1, 2, 3 i stedet for 1.1, 1.2. Passer en handout
% som handler om én ting.
%
% Med [blokkavsnitt] skilles avsnittene med luft (en halv
% grunnlinje) i stedet for innrykk. Valgene kan kombineres.
% ============================================================
\usepackage{handoutstyle}
% \usepackage[unummerert]{handoutstyle}
% \usepackage[blokkavsnitt]{handoutstyle}
% \usepackage[unummerert,blokkavsnitt]{handoutstyle}


% ============================================================
% VALGFRIE HANDOUT-JUSTERINGER  (fjern % for å slå på)
% ============================================================

% --- Av med sidetall (trenger BEGGE linjer) ---
% \maketitle tvinger "plain" på side 1, så fjern også %-tegnet på
% \thispagestyle{empty} etter \maketitle for å fjerne tallet der.
% \pagestyle{empty}


% ============================================================
% EGENDEFINERTE PAKKER
%
% Legg til prosjektspesifikke pakker her.
% Ikke sett pakker andre steder i filen.
% ============================================================

% --- Matematikk (standard) ---
\usepackage{amsmath, amssymb}

% --- Figurer ---
\usepackage{graphicx}                % figurer
\graphicspath{{figurer/}}            % figurer ligger i figurer/

% --- Valgfritt ---
% \usepackage{physics}              % \dd, \ket, \bra, ...
\usepackage{tikz}                   % native TikZ-figurer


% ============================================================
% EGENDEFINERTE KOMMANDOER
%
% Prosjektspesifikke makroer.
%
% \separator-linjen leveres av handoutstyle.sty.
% ============================================================

% \newcommand{\dd}{\mathrm{d}}
% \newcommand{\Ham}{\mathcal{H}}


% ============================================================
% METADATA
% ============================================================
\title{Handouttittel}
\date{Våren 2026 $\cdot$ Kursnavn}


% ============================================================
% DOKUMENT
% ============================================================
\begin{document}

\maketitle
% \thispagestyle{empty}   % side 1: fjerner tallet (se \pagestyle over)


% ============================================================
% SEKSJON
% ============================================================
\section{Energi og bevaring}

En partikkel med masse $m$ beveger seg i et endimensjonalt
potensial $V(x)$. Den totale mekaniske energien er
\begin{equation}
  E = \tfrac{1}{2} m v^2 + V(x),
\end{equation}
som er bevart når ingen ytre krefter virker.

% viktig-boksen løfter en sentral definisjon fra prosaen;
% nummereringen fortsetter i samme rekke som utenfor.
\begin{viktig}
\begin{definisjon}
En kraft er \emph{konservativ} dersom den kan skrives som den
negative gradienten av et potensial, $F(x) = -\,\mathrm{d}V/\mathrm{d}x$.
\end{definisjon}
\end{viktig}

% hovedresultat rammer inn dokumentets sentrale resultat;
% tittelargumentet kan sløyfes for en ren ramme.
\begin{hovedresultat}[Energibevaring]
\begin{teorem}
For en konservativ kraft er den totale mekaniske energien $E$
konstant langs enhver løsning av bevegelsesligningen.
\end{teorem}
\end{hovedresultat}

\begin{merknad}
Det omvendte holder ikke generelt: dissipative systemer har
ikke noe slikt potensial, og $E$ er ikke bevart.
\end{merknad}

\subsection{Potensialenergidiagrammer}

En graf av $V(x)$ mot $x$ oppsummerer den kvalitative bevegelsen
med ett blikk: partikkelen er begrenset til områder der
$E \ge V(x)$, og vendepunkter opptrer der $E = V(x)$,
som illustrert i figur~\ref{fig:potensial}.

% TikZ-figureksempel: bruk darkorange til energinivåer / strukturelementer,
% darkolive til kurver (darkblue og darkplum til flere serier), svart
% til akser, grå til sekundære etiketter. darkred er reservert for
% separatorlinjen -- bruk den ikke i figurer.
\begin{figure}[ht]
  \centering
  \begin{tikzpicture}[>=stealth]
    % Akser
    \draw[->,thick] (-2.2,0) -- (2.2,0) node[right,font=\small] {$x$};
    \draw[->,thick] (0,-0.2) -- (0,3.5) node[above,font=\small] {$V$};
    % Parabel V(x) ~ x^2 (darkolive)
    \draw[darkolive,thick,domain=-1.85:1.85,samples=60]
      plot (\x, {\x*\x});
    \node[darkolive,right,font=\small] at (1.7,3.1) {$V(x)$};
    % Energinivå (darkorange)
    \draw[darkorange,thick] (-1.6,2.56) -- (1.6,2.56);
    \node[right,font=\small,darkorange] at (1.65,2.56) {$E$};
    % Klassiske vendepunkter (grå stiplet)
    \draw[gray,thin,dashed] (-1.6,0) -- (-1.6,2.56);
    \draw[gray,thin,dashed] ( 1.6,0) -- ( 1.6,2.56);
    \node[below,font=\small,gray] at (-1.6,0) {$-x_0$};
    \node[below,font=\small,gray] at ( 1.6,0) {$x_0$};
  \end{tikzpicture}
  \caption{Harmonisk potensial (\textcolor{darkolive}{olivengrønn}) med
           energinivå $E$ (\textcolor{darkorange}{oransje}) og klassiske
           vendepunkter $\pm x_0$ (grå).}
  \label{fig:potensial}
\end{figure}

Underseksjonsoverskriften over ligger ett størrelsessteg under
seksjonstittelen, men holder seg på kroppsstørrelse --- nok til
å markere et finere nivå når handout-en brukes som
forelesningsnotater.


% ============================================================
% OPPGAVE
% ============================================================
\begin{oppgave}[Harmonisk oscillator]
Betrakt en partikkel med masse $m$ i potensialet
$V(x) = \tfrac{1}{2}kx^2$.

\begin{deloppgaver}
  \item Skriv opp bevegelsesligningen ved hjelp av Newtons
        andre lov.
  \item Vis at
        $E = \tfrac{1}{2}mv^2 + \tfrac{1}{2}kx^2$
        er bevart langs løsningene.
\end{deloppgaver}

Resultatet kan verifiseres ved å derivere $E$ med hensyn på
tid og bruke bevegelsesligningen fra deloppgave~(a).

\begin{deloppgaver}[resume]
  \item Finn vinkelfrekvensen for små svingninger.
\end{deloppgaver}
\end{oppgave}

\separator


% ============================================================
% OPPGAVE
% ============================================================
\begin{oppgave}[Skråplan]
En kloss med masse $m$ glir nedover et friksjonsfritt
skråplan som danner vinkelen $\theta$ med horisontalen.

\begin{deloppgaver}
  \item Finn akselerasjonen langs skråplanet.
  \item Hvor lang tid tar det å gli en strekning $L$ fra ro?
  \item Hva er farten ved bunnen?
\end{deloppgaver}
\end{oppgave}


% ============================================================
% SEKSJON
% ============================================================
\section{Sirkelbevegelse}

\begin{figure}[t]
  \centering
  \includegraphics[width=0.5\linewidth]{Energy_levels}
  \caption{Skjematisk figur for temaet.}
  \label{fig:sirkel}
\end{figure}

Ved jevn sirkelbevegelse har sentripetalakselerasjonen
størrelsen $a_s = v^2/r$, rettet inn mot sentrum av sirkelen.


% ============================================================
% OPPGAVE
% ============================================================
\begin{oppgave}[Sirkulær bane]
\label{opg:bane}
En satellitt går i en sirkulær bane med radius $r$ rundt en
planet med masse $M$.

\begin{deloppgaver}
  \item Vis at banehastigheten er $v = \sqrt{GM/r}$, hvor
        $G$ er Newtons gravitasjonskonstant.
  \item Uttrykk omløpstiden $T$ ved $r$, $M$ og $G$.
\end{deloppgaver}
\end{oppgave}


% ============================================================
% SEKSJON
% ============================================================
\section{Numerisk eksempel}

\verb|python|-omgivelsen setter kode med Gruvbox Light
syntaksutheving. Det korte skriptet nedenfor verifiserer
banehastigheten og omløpstiden fra oppgave~\ref{opg:bane} numerisk for
tre representative radier.

\begin{python}
import numpy as np

G = 6.674e-11   # gravitasjonskonstant  (m^3 kg^-1 s^-2)
M = 5.972e24    # jordmasse             (kg)

baner = {"Lav bane":   6.40e6,
         "ISS (ca.)":  6.78e6,
         "GEO":        4.22e7}

for navn, r in baner.items():
    v = np.sqrt(G * M / r)           # banehastighet  [m/s]
    T = 2 * np.pi * r / v / 3600    # omløpstid      [timer]
    print(f"{navn:12s}  v = {v/1e3:.3f} km/s   T = {T:.2f} t")
\end{python}


\end{document}
